\section{Capability Registry}

The capability registry is the default authority substrate. A capability is normally a live grant in registry state, not a token passed into every method.

\subsection{Registry grants}

Example registry records:

\begin{lstlisting}
(RoleGrant
  (contract TokenA)
  (role Minter)
  (holder BridgeA)
  (admin Governance)
  (active True))

(AllowanceGrant
  (token TokenA)
  (owner Alice)
  (spender Dex)
  (asset USDC)
  (limit 1000)
  (spent 250)
  (expiry Block-900000)
  (revoked False))

(OracleGrant
  (oracle PriceOracleA)
  (feed ETH-USD)
  (holder OracleBot7)
  (right update-price)
  (active True))
\end{lstlisting}

\subsection{Registry-first authority}

The normal authorization path is:

\begin{lstlisting}
method call
  -> policy lookup
  -> derive grant key from method arguments and caller
  -> live registry lookup
  -> authorize or reject
\end{lstlisting}

For example, \code{transferFrom(Alice, Bob, USDC, 10)} called by \code{Dex} derives the canonical lookup:

\begin{lstlisting}
AllowanceGrant(TokenA, owner=Alice, spender=Dex, asset=USDC)
\end{lstlisting}

The method succeeds only if that grant exists, is active, is not expired or revoked, and has sufficient remaining amount.

\subsection{Registry requirements}

\req{SCS-REG-001 --- Registry-first authority}{The default capability certification mechanism \MUST{} be live lookup in an authoritative capability registry.}

\req{SCS-REG-002 --- Contract-scoped registries}{Each contract \MAY{} own its own capability registry. A runtime \MAY{} also provide a global registry, but authority records \MUST{} be explicitly scoped by contract and resource.}

\req{SCS-REG-003 --- Guarded registry mutation}{Only authorized methods \MAY{} create, update, consume, revoke, or delete registry grants.}

\req{SCS-REG-004 --- Use-time lookup}{Protected operations \MUST{} check the registry at use time.}

\req{SCS-REG-005 --- No authority from syntax}{A \MeTTa{} term that merely resembles a grant \MUSTNOT{} authorize anything unless it exists in the authoritative registry or is accepted by an approved certificate adapter.}

\req{SCS-REG-006 --- Canonical lookup keys}{Registries \SHOULD{} use canonical keys.}

Examples:

\begin{lstlisting}
Role:
    (contract, role, holder)

Allowance:
    (token, owner, spender, asset)

Oracle:
    (oracle, feed, updater)

Governance:
    (contract, action, holder)
\end{lstlisting}

\req{SCS-REG-007 --- Atomic consumption}{If a grant use consumes allowance, quota, nonce, or spend limit, that consumption \MUST{} be staged and committed atomically with the protected operation.}

\req{SCS-REG-008 --- Revocation visibility}{A revoked grant \MUST{} fail once the revocation is visible in the current execution state.}

\req{SCS-REG-009 --- No stale authorization cache}{A cached authorization result \MUSTNOT{} authorize a later operation after the underlying grant has changed, unless snapshot semantics explicitly permit it.}

\req{SCS-REG-010 --- Deterministic lookup}{Registry lookup \MUST{} be deterministic under consensus execution.}

\req{SCS-REG-011 --- Registry schema}{Capability registries \SHOULD{} define schemas for each grant class.}

\req{SCS-REG-012 --- Registry invariants}{The registry \MUST{} enforce grant-class invariants such as nonnegative limits, valid expiry domains, valid holder identifiers, valid resource identifiers, and valid right names.}

\subsection{Registry consumption}

For spend-like grants, authorization has a side effect. For \code{transferFrom}, the runtime must stage:

\begin{lstlisting}
AllowanceGrant.spent := AllowanceGrant.spent + amount
\end{lstlisting}

That staged update \MUST{} commit only if the corresponding balance transfer also commits.
